\ficha{Bordo e Rockoff (1996), Good Housekeeping Seal}{bordorockoff1996}

\begin{identificacao}
BORDO, Michael D.; ROCKOFF, Hugh. The gold standard as a ``good housekeeping
seal of approval''. \emph{The Journal of Economic History}, v. 56, n. 2,
p. 389-428, 1996.

Artigo de história econômica quantitativa, em inglês, lido integralmente na
versão de working paper: NBER Working Paper 5340, novembro de 1995, 37 páginas.
As páginas citadas abaixo são as do working paper e não coincidem com a
paginação do artigo publicado.
\end{identificacao}

\bloco{Problema e tese}
Por que países periféricos aceitaram o custo de manter, ou de tentar manter, a
conversibilidade em ouro entre 1870 e 1914? A resposta é que a adesão
funcionava como sinal de probidade financeira e barateava o crédito nos centros
europeus. Isso explicaria a determinação de importadores de capital em ficar no
ouro, e oferece alternativa às explicações que apelam ao conflito interno entre
credores e devedores ou à preferência irracional pelo metal.

\bloco{Argumento}
\begin{enumerate}[leftmargin=*,nosep]
\item O padrão-ouro é lido como regra contingente, ou regra com cláusula de
escape: a autoridade mantém o preço fixo do ouro exceto em emergência bem
compreendida, tipicamente uma guerra grande, com o entendimento público de que
a suspensão dura o tempo da emergência e que a retomada será na paridade
original (p. 4).
\item A regra tem de ser transparente e conter poucas contingências. Invocar a
escape em crise financeira ou em choque de termos de troca é mais suspeito,
porque a fonte da contingência é difícil de verificar (p. 5). São exemplos de
discrição, isto é, de quebra da regra,
\chave{postponement of resumption after the war and reasonable delay period had
passed, and pegging to specie at a devalued parity}{p. 5}
\item Daí a assimetria. Os países centrais suspenderam em guerra e voltaram à
paridade anterior. Os latino-americanos suspenderam também por crise financeira
e por choque de termos de troca, e retornaram em paridade depreciada, com maior
crescimento monetário, déficit e inflação:
\chave{For them gold convertibility was the exception rather than the rule}{p. 6}
\item O enforcement examinado é o acesso a capital: se aderir sinaliza boa
conduta, quem sempre aderiu paga juros menores em Londres (p. 7). Os autores
admitem que só mediram o preço do crédito, não o volume atraído (p. 7).
\item O resultado é duplo. A \emph{dummy} de estar em ouro é negativa e
significativa na amostra de títulos-ouro, e os autores resumem o achado em 40
pontos-base, o ``haircut'' cobrado por não estar em ouro (p. 17). O efeito maior
passa pelos betas, que medem compromisso de longo prazo: Argentina, Brasil e
Chile pagaram de dois a três pontos percentuais acima do consol britânico,
Canadá, Austrália e Estados Unidos pagaram cerca de um ponto, e a Itália, de
adesão formal pior mas política ortodoxa, pouco mais que eles (p. 18).
\item A conclusão registra o custo da regra: quem aderiu abriu mão da
flexibilidade de responder a choques adversos de oferta por política financeira
expansionista e alteração do câmbio (p. 18).
\end{enumerate}

\bloco{Conceitos e definições operacionais}
Regra é mecanismo de compromisso crível contra a inconsistência temporal, não
automatismo impessoal: impede o governo de emitir moeda fiduciária por surpresa
para capturar senhoriagem ou de dar calote na dívida (p. 3). Credibilidade é
sinal em mundo de informação assimétrica, pois a adesão informa ao credor que
custos o devedor suportará para não dar calote (p. 2). Adesão é operacionalizada
como variável binária por país-ano (p. 13); compromisso de longo prazo, pelo
beta, coeficiente do diferencial médio da amostra, análogo do risco sistemático
do CAPM (p. 13). Sombrear o ouro é conduzir política fiscal e monetária
compatível com a paridade sem conversibilidade formal, caso de Espanha,
Portugal e da Itália posterior a 1903 (p. 8). Títulos-ouro isolam o risco-país
do risco cambial, porque prometem juros e principal em ouro (nota 2, p. 2).

\bloco{Evidência e método}
Observações anuais de 1870 a 1914 para nove países periféricos: rendimento de
títulos governamentais de longo prazo cotados em Londres, os italianos em Paris,
mais câmbio, renda real, déficit fiscal e oferta de moeda. A amostra é partida
entre títulos-ouro, sete países, e títulos em papel, cinco, e os países se
distribuem em quatro grupos, dos sempre aderentes (Canadá e Austrália) aos que
quebraram a regra suspendendo de forma intermitente e desvalorizando (Argentina,
Brasil e Chile) (p. 8). A cronologia de adesão da Tabela 1 (p. 33) vem de Bordo
e Schwartz (1994). Os pesos da média ponderada são a participação de cada país
na dívida detida pela Grã-Bretanha em 1914, segundo Feis, com o Brasil em 6,6\%
(nota 18, p. 13). A estimação vai de mínimos quadrados por país, com correção de
autocorrelação por Cochrane-Orcutt, a um modelo SUR agrupado que restringe o
coeficiente da \emph{dummy} e o dos fundamentos a serem iguais entre países
(p. 15). Não há contrafactual formal; a comparação é entre grupos de adesão.

\bloco{O Brasil na amostra}
O Brasil está na amostra, no grupo dos que quebraram a regra. A narrativa
brasileira ocupa um parágrafo (p. 9), com Peláez e Suzigan (1976) e Fritsch e
Franco (1992) como fontes, e vai do bimetalismo de 1808 à revolução republicana
de 1889, que coincide com o fim da conversibilidade. Sobre 1906 o texto diz que
o Brasil restaurou a conversibilidade para impedir a valorização do mil-réis,
prejudicial aos exportadores de café e borracha, e criou uma ``Conversion
Office'' com limite de emissão a paridade nova (p. 9). A Tabela 1 (p. 33)
registra 1906-1914 como ouro, motivo da mudança guerra, mudança de paridade sim.
A série de juros brasileira começa em 1892 (p. 12). Nas regressões agrupadas o
beta brasileiro é 1,49 na média não ponderada e 2,06 na ponderada, o mais alto
da amostra ao lado do argentino (Tabela 4, p. 36). O texto não menciona o
Convênio de Taubaté, nem a taxa de 15 pence, nem o par legal de 27 pence, nem
qualquer agente brasileiro.

\bloco{Agentes e vertentes}
Filiação teórica: Kydland e Prescott e Barro e Gordon para regra contra
discrição; Bordo e Kydland (1995) e Bordo e Schwartz (1994) para a regra
contingente; Stiglitz e Weiss para racionamento de crédito sob informação
assimétrica; Sharpe e Lintner para o CAPM. A nota 3 (p. 3) situa a
\emph{currency school} e a \emph{banking school} como a disputa inglesa entre
emissão automática contra reserva metálica e discrição dos diretores do Banco da
Inglaterra. Na historiografia aparecem Eichengreen, Officer, Giovannini, Cortés
Conde para a Argentina, Llona-Rodríguez para o Chile, Martín Aceña para a
Espanha, Reis e Macedo para Portugal.

\bloco{Críticas e limites}
Os autores medem preço do crédito, não volume (p. 7). A convergência geral dos
rendimentos com o consol depois de 1900 confunde o efeito da adesão, e o texto
reconhece isso justamente no caso relevante aqui: as taxas de Argentina e Brasil
caem depois de 1900, quando ambos estavam em ouro, mas é difícil determinar se a
queda veio da adesão ou de outro fator (p. 12). Na regressão individual do
Brasil, com 24 observações, a \emph{dummy} de ouro é insignificante e o beta
perde significância na especificação com variáveis de política (Tabela 2,
p. 34); o resultado brasileiro depende da restrição de coeficiente comum na SUR.
A \emph{dummy} binária não distingue retorno na paridade original de retorno em
paridade desvalorizada, embora a seção conceitual trate a segunda como quebra
(p. 5): o Brasil de 1906 conta como país em ouro apesar de a Tabela 1 registrar
mudança de paridade. O R\textsuperscript{2} simulado da SUR tem propriedades
estatísticas desconhecidas, e os autores o dizem (p. 15). Obstfeld e Taylor
(2003) e Ferguson e Schularick (2008), também fichados neste caderno, disputam o
peso do ouro frente a fundamentos fiscais e a vínculo imperial.

\begin{dialogo}
\item Procurar, nos editoriais de 1906 e nos Documentos Parlamentares, o
argumento do crédito externo como razão para a conversibilidade: crédito,
confiança, praça de Londres, juro dos empréstimos. Testa se a hipótese do selo
circulava como razão pública brasileira ou se a defesa da Caixa de Conversão se
fez apenas pela estabilidade cambial do exportador.
\item Procurar como se tratou a distância entre a paridade da Caixa e o par
legal de 27 pence. Bordo e Rockoff classificam fixar paridade desvalorizada como
quebra da regra (p. 5), e a imprensa ortodoxa pode ter usado argumento
equivalente contra David Campista, ou pode ter separado estabilidade de paridade
nominal.
\item Procurar menção ao precedente argentino de 1899 e à Caja de Conversión
como caso de queda do prêmio de risco. Testa se a comparação platina entrou no
debate com conteúdo financeiro ou só como analogia institucional.
\item Procurar, em 1914, se a suspensão do troco foi justificada com linguagem
de emergência de guerra, que é a cláusula de escape legítima do modelo, ou como
abandono da paridade.
\end{dialogo}

\usonocapitulo{Parte I. Sustenta que a adesão ao padrão-ouro tinha preço
mensurável no mercado de capitais e que o prêmio pago pelos periféricos
respondia menos ao estar em ouro num dado ano do que à reputação acumulada, o
que autoriza tratar credibilidade como estoque e não como decisão pontual. Serve
também para situar o Brasil de 1906 na fronteira do argumento: o país entra na
amostra como quebrador da regra e sai dela com o beta mais alto ao lado do
argentino, embora a criação da Caixa de Conversão apareça em duas frases e sem
nenhum agente nomeado. O parágrafo final sobre currency boards contemporâneos
(p. 19) faz a ligação com a Parte IV.}
